% Part: sets-functions-relations
% Chapter: infinite
% Section: dedekinds-proof
%
\documentclass[../../../include/open-logic-section]{subfiles}

\begin{document}

\olfileid{sfr}{infinite}{dedekindsproof}

\olsection[ডেডেকিন্ডের ``প্রমাণ'']{অসীম সেটের অস্তিত্ব নিয়ে ডেডেকিন্ডের ``প্রমাণ''}

এই অধ্যায়ে আমরা ডেডেকিন্ড বীজগঠনের সাহায্যে স্বাভাবিক সংখ্যাগুলির একটি সেটতাত্ত্বিক আলোচনা দিয়েছি। \olref[arith][ref]{sec}-এ আমরা বিশ্লেষণের সংখ্যাতাত্ত্বিক নির্মাণের দার্শনিক তাৎপর্য (অন্যান্য বিষয়ের সঙ্গে) বিবেচনা করেছিলাম। এবার এখানে আমরা কী অর্জন করেছি, তার তাৎপর্য বিবেচনা করা উচিত।

\olref[sfr][arith][]{chap} জুড়ে আমরা স্বাভাবিক সংখ্যাগুলিকে প্রদত্ত ধরে নিয়ে তাদের সাহায্যে পূর্ণসংখ্যা, মূলদ সংখ্যা ও বাস্তব সংখ্যা স্পষ্টভাবে নির্মাণ করেছি। এই অধ্যায়ে আমরা স্বাভাবিক সংখ্যাগুলির কোনো স্পষ্ট নির্মাণ দিইনি। শুধু দেখিয়েছি যে \emph{যেকোনো ডেডেকিন্ড-অসীম সেট দেওয়া থাকলে} এমন একটি সেট সংজ্ঞায়িত করা যায়, যা আমরা~$\Nat$-এর কাছে যে আচরণ চাই, ঠিক তেমনই আচরণ করবে।

তাই \emph{কোন বস্তুগুলি স্বাভাবিক সংখ্যা}—এমন কোনো অধিবিদ্যাগত প্রশ্নের উত্তর দিয়েছি বলে আমরা নিশ্চয়ই দাবি করতে পারি না। আর সেটিই ভালো। কারণ \olref[sfr][arith][ref]{sec}-এ আমরা জোর দিয়ে বলেছিলাম, ডেডেকিন্ড কর্তনের সেট হিসেবে $\Real$-এর সংজ্ঞাকে সুবিধাজনক একটি নির্ধারণ না ভেবে \emph{আবিষ্কার} ভাবলে ভুল হবে। মুখ্য পর্যবেক্ষণ হলো, ডেডেকিন্ড কর্তনগুলি বাস্তব সংখ্যার প্রধান গাণিতিক ধর্মগুলি দৃষ্টান্তায়িত করে। এখানেও একই কথা: \emph{যেকোনো} ডেডেকিন্ড বীজগঠন স্বাভাবিক সংখ্যার প্রধান গাণিতিক ধর্মগুলি দৃষ্টান্তায়িত করে। (বস্তুত, ডেডেকিন্ড সব ডেডেকিন্ড বীজগঠন যে \emph{সমরূপ}, তা প্রমাণ করে এই বক্তব্যকে সুপ্রতিষ্ঠিত করেছিলেন (\citeyear[Theorems
132--3]{Dedekind1888})। তাই অনেক সমকালীন ``গঠনবাদী'' যে ডেডেকিন্ডকে অগ্রদূত হিসেবে উল্লেখ করেন, তাতে আশ্চর্যের কিছু নেই।)

তাছাড়া, স্বাভাবিক সংখ্যার তত্ত্বকে কীভাবে একটি অনানুষ্ঠানিক সরল সেটতত্ত্বের মধ্যে অন্তর্ভুক্ত করা যায়, তা আমরা দেখিয়েছি। সেই সেটতত্ত্ব এখনও বেশ অনানুষ্ঠানিক, কিন্তু তাতে (আপাতদৃষ্টিতে) স্বাভাবিক সংখ্যাগুলিকে প্রদত্ত ধরে নেওয়া হয় না। তাই ডেডেকিন্ডের উচ্চাভিলাষী নিজস্ব প্রকল্পটি বাস্তবায়নের পথে হয়তো আমরা এগোচ্ছি। তিনি প্রকল্পটি এভাবে ব্যাখ্যা করেছিলেন:
\begin{quote}
	বিজ্ঞানে যা প্রমাণযোগ্য, তা প্রমাণ ছাড়া বিশ্বাস করা উচিত নয়। এই দাবি যুক্তিসঙ্গত মনে হলেও আমার মতে, সর্বাধিক সরল বিজ্ঞানের ভিত্তি স্থাপনের এমনকি সর্বাধুনিক পদ্ধতিগুলিতেও এটি পূরণ হয়নি; অর্থাৎ সংখ্যার তত্ত্ব নিয়ে কাজ করে যুক্তিবিদ্যার যে অংশ, সেখানেও নয়। পাটীগণিতকে (বীজগণিত, বিশ্লেষণ) যুক্তিবিদ্যার নিছক একটি অংশ বলার সময় আমি বোঝাতে চাই যে সংখ্যা-ধারণাকে আমি স্থান ও কালের ধারণা বা স্বজ্ঞা থেকে সম্পূর্ণ স্বাধীন বলে মনে করি—বরং তাকে চিন্তার বিশুদ্ধ বিধিগুলির প্রত্যক্ষ ফল বলে মনে করি।
	\citep[preface]{Dedekind1888}
\end{quote}
ডেডেকিন্ডের সাহসী ধারণাটি হলো এই। শুধু (অনানুষ্ঠানিক) সেটতত্ত্ব ব্যবহার করে কীভাবে স্বাভাবিক সংখ্যা নির্মাণ করা যায়, তা আমরা এইমাত্র দেখিয়েছি। \olref[sfr][arith][]{chap}-এ দেখেছি, স্বাভাবিক সংখ্যা ও কিছু সেটতত্ত্ব দেওয়া থাকলে কীভাবে বাস্তব সংখ্যা নির্মাণ করা যায়। তাই হয়তো ``পাটীগণিত (বীজগণিত, বিশ্লেষণ)'' শেষ পর্যন্ত ``যুক্তিবিদ্যার নিছক একটি অংশ'' হয়ে ওঠে (ডেডেকিন্ডের ``যুক্তিবিদ্যা'' শব্দটির প্রসারিত অর্থে)।

ধারণাটি এই। কিন্তু একটু থামো। আমাদের ডেডেকিন্ড বীজগঠনের নির্মাণ (যা স্বাভাবিক সংখ্যার প্রতিস্থাপক) একটি ডেডেকিন্ড-অসীম সেটের অস্তিত্বের উপর শর্তাধীন। (\olref[sfr][infinite][dedekind]{thm:DedekindInfiniteAlgebra}-এ ফিরে তাকাও।) ``যুক্তিবিদ্যা'' বা ``চিন্তার বিশুদ্ধ বিধি'' দিয়ে যদি ডেডেকিন্ড-অসীম সেটের অস্তিত্ব প্রতিষ্ঠা করা না যায়, তবে প্রকল্পটি থেমে যায়।

তাহলে ``চিন্তার বিশুদ্ধ বিধি'' দিয়ে কি ডেডেকিন্ড-অসীম সেটের অস্তিত্ব প্রতিষ্ঠা করা \emph{যায়}? ডেডেকিন্ডের প্রচেষ্টা ছিল এই:
\begin{quote}
  আমার নিজের চিন্তার জগৎ, অর্থাৎ আমার চিন্তার বিষয় হতে পারে এমন সমস্ত বস্তুর সমগ্রতা $S$, অসীম। কারণ $s$ যদি~$S$-এর একটি উপাদান নির্দেশ করে, তবে~$s$ আমার চিন্তার বিষয় হতে পারে—এই চিন্তা $s'$ নিজেই~$S$-এর একটি উপাদান। একে যদি উপাদান~$s$-এর একটি প্রতিবিম্ব $\phi(s)$ হিসেবে ধরি, তবে \dots~$S$ [ডেডেকিন্ড-]অসীম, যা প্রমাণ করার ছিল।
	\citep[\S66]{Dedekind1888}
\end{quote}
মূলত অত্যন্ত কঠোর গাণিতিক প্রমাণে ভরা একটি বইয়ের মাঝখানে এমন কথা পাওয়া বেশ বিস্ময়কর। দুটি মন্তব্য করা দরকার।

প্রথমত, এই ``প্রমাণ''-এ এমন কোনো চরিত্র প্রায় নেই, যাকে আমরা এখন ``গাণিতিক'' বলে চিনব। এতে মনস্তাত্ত্বিক বস্তু (চিন্তা), তাও শুধু \emph{সম্ভাব্য} বস্তুর কথা বলা হয়েছে।

দ্বিতীয়ত, অন্তত আমরা যেভাবে ডেডেকিন্ড বীজগঠন উপস্থাপন করেছি, তাতে এই ``প্রমাণ''-এর একটি সরল কারিগরি ঘাটতি আছে। ডেডেকিন্ডের যুক্তি সফল হলেও তা শুধু প্রতিষ্ঠা করে যে অসীমসংখ্যক বস্তু আছে (বিশেষত, অসীমসংখ্যক চিন্তা)। কিন্তু~$S$-কে অসীমসংখ্যক !!{element}s নিয়ে গঠিত একটি একক \emph{সেট} হিসেবে দেখার, বহুবচনে~$S$-কে \emph{কয়েকটি বস্তু} না ভাবার, কারণও ডেডেকিন্ডকে দিতে হবে।

ডেডেকিন্ড এখানে কোনো ফাঁক দেখতে পাননি—এই তথ্যটি হয়তো ইঙ্গিত করে যে তাঁর ``সমগ্রতা'' শব্দের ব্যবহার আমাদের ``সেট'' শব্দের ব্যবহারের সঙ্গে হুবহু মেলে না।\footnote{বস্তুত, তা মেলেনি ভাবার অন্য কারণও আমাদের আছে; \citet[p.~23]{Potter2004} দেখো।} তবে তাতে খুব আশ্চর্য হওয়ার কিছু নেই। গত দুই অধ্যায়ে আমরা যে প্রকল্প অনুসরণ করেছি—স্বাভাবিক সংখ্যার একটি ``নির্মাণ'', এবং তা থেকে পূর্ণসংখ্যা, বাস্তব সংখ্যা ও মূলদ সংখ্যার ``নির্মাণ''—তার সবটাই অনানুষ্ঠানিকভাবে করা হয়েছে। ঠিক কোন সেট কখন বিদ্যমান, সে বিষয়ে বেশি কোনো সাধারণ নীতি নির্ধারণ না করেই প্রয়োজনমতো কখনও এই সেট, কখনও ওই সেট নিয়ে নিয়েছি। কিন্তু আমরা জানি, \emph{কিছু} সাধারণ নীতি আমাদের দরকার; নইলে রাসেলের কূটাভাসে পড়ব।

এবার আমাদের এই অনানুষ্ঠানিকতার সীমা অতিক্রম করার সময় হয়েছে।

\end{document}
